\section*{Chapitre \ref{ch:b2:guided-waves} --- Ondes guidées et cavités}

\begin{solution}{exo:b2:guided-waves:1}
$f_n = nc/2a = 7.5$, $15$, \qty{22.5}{GHz}. À \qty{12}{GHz} seul $n = 1$~: $k_g =
(2\pi/c)\sqrt{f^2 - f_1^2} = \qty{196}{rad/m}$, $\lambda_g = \qty{3.2}{cm}$, $v_\varphi = c/0.78
= \qty{3.8e8}{m/s}$, $v_g = \qty{2.3e8}{m/s}$. À \qty{6}{GHz}~: $\ell = c/2\pi\sqrt{f_1^2 -
f^2} = \qty{1.1}{cm}$.
\end{solution}

\begin{solution}{exo:b2:guided-waves:2}
TE$_{10}$ \qty{6.55}{GHz}, TE$_{20}$ \qty{13.1}{GHz}, TE$_{01}$ \qty{14.7}{GHz}~: monomode
de $6.55$ à \qty{13.1}{GHz}. À \qty{10}{GHz}~: $\sqrt{1 - (6.55/10)^2} =
0.756$, $\lambda_g = \qty{4.0}{cm}$, $v_g = 0.76c$, $\cos\theta = 0.756$, $\theta = \qty{41}{\degree}$.
$b \approx a/2$ maintient TE$_{01}$ au-dessus de TE$_{20}$ (la plus large bande monomode)
tout en laissant le plus grand écart au champ (puissance avant claquage).
\end{solution}

\begin{solution}{exo:b2:guided-waves:3}
$f = 1.5 \times 10^8\sqrt{(m/0.3)^2 + (n/0.3)^2 + (p/0.2)^2}$~: \qty{0.71}{GHz}, \qty{0.90}{GHz},
\qty{1.35}{GHz}, \qty{1.95}{GHz}. $N = 8\pi Vf^3/3c^3 = 8\pi \times 0.018 \times 1.56 \times 10^{28}/
8.1 \times 10^{25} \approx 260$. Cube~: $f = c\sqrt2/2L$, $L = \qty{2.3}{cm}$.
\end{solution}

\begin{solution}{exo:b2:guided-waves:4}
$\mathrm{NA} = \sqrt{1.48^2 - 1.46^2} = 0.24$, \qty{14}{\degree}~; $\Delta t = (n_1L/c)(n_1/n_2
- 1) = \qty{68}{ns/km}$~; \qty{135}{ns} sur \qty{2}{km}~: \qty{7}{Mbit/s}~; avec $0.003$~:
\qty{10}{ns/km}, \qty{50}{Mbit/s}.
\end{solution}

\begin{solution}{exo:b2:guided-waves:5}
(a) $\underline B_y = (k_g/\omega)\underline E$, $\underline B_z = (\iu/\omega)\partial_y\underline E = (\iu\pi/a\omega)
E_0\cos(\pi y/a)\eu^{\iu(\omega t - k_gz)}$~: en quadrature. (b) En $y = 0$, $B_y = 0$ et
$\vect j_s = \vect e_y\wedge\vect B/\mu_0 = (B_z/\mu_0)\vect e_x$~: selon le champ. (c)
$\langle\Pi_z\rangle = E_0^2k_g\sin^2(\pi y/a)/2\mu_0\omega$, puissance par unité de largeur
$E_0^2k_ga/4\mu_0\omega$~; $\langle\Pi_y\rangle \propto \operatorname{Re}(\underline E\,\iu\underline B_z^*)
= 0$. (d) Énergie moyenne par unité de longueur et de largeur $\tfrac14\varepsilon_0E_0^2a$ (les moitiés électrique
et magnétique valant chacune $\tfrac18\varepsilon_0E_0^2a$)~; le rapport vaut $c^2k_g/\omega = v_g$.
\end{solution}

\begin{solution}{exo:b2:guided-waves:6}
(a) $f_1 = \qty{150}{GHz} \gg f$~: $\ell \approx a/\pi = \qty{0.32}{mm}$~; $\eu^{-0.5/0.32}$~:
\qty{-14}{dB} en amplitude, $-27$ en puissance. (b) $f_1 = \qty{750}{MHz}$~; à \qty{100}{MHz}
$\ell = \qty{6.4}{cm}$~: \qty{-136}{dB} sur un mètre. (c) $f_1 = \qty{15}{GHz}$, $\ell =
\qty{3.2}{mm}$~: \qty{27}{dB/cm}. (d) L'onde n'y tient pas~; l'évanescence est
géométrique, non dissipative.
\end{solution}

\begin{solution}{exo:b2:guided-waves:7}
(a) $P = \tfrac14\varepsilon_0cE_0^2ab\sqrt{1 - f_{10}^2/f^2} = \qty{120}{kW}$. (b) $\times9$~:
\qty{1}{MW}~; l'azote sous pression ou le SF$_6$ élèvent le champ de claquage. (c)
$j_s \approx E_0/\mu_0c = \qty{2.7}{kA/m}$~; $\tfrac12R_sj_s^2 = \qty{9e4}{W/m^2}$ sur $2a =
\qty{0.046}{m^2}$ par mètre~: \qty{4}{kW/m}, soit $3.5\%$ par mètre, environ \qty{0.15}{dB/m}.
(d) Trois fois moins que le câble, et des mégawatts au lieu de
kilowatts.
\end{solution}

\begin{solution}{exo:b2:guided-waves:8}
(a) $f = c\sqrt2/2L = \qty{2.12}{GHz}$. (b) $\delta = \qty{1.4}{\micro m}$, $R_s = \qty{12}{m\ohm}$~;
$W = \varepsilon_0E_0^2V/8 = \qty{1.1e-15}{}E_0^2$~; $P = \tfrac12R_s(E_0/377)^2 \times 6L^2 =
\qty{2.5e-9}{}E_0^2$~: $Q = \omega W/P \approx 6000$. (c) \qty{360}{kHz}~; \qty{0.45}{\micro s}. (d)
$Q \sim 10^{10}$~: un champ de dizaines de mégavolts par mètre est entretenu par
des kilowatts au lieu de gigawatts.
\end{solution}

\begin{solution}{exo:b2:guided-waves:9}
(a) Les points $(m, n, p)$ vérifiant $m^2 + n^2 + p^2 \le (2fL/c)^2$ remplissent un huitième
de sphère~: $\tfrac18\cdot\tfrac43\pi(2fL/c)^3$, multiplié par deux polarisations~:
$8\pi L^3f^3/3c^3$. (b) $8\pi f^2/c^3$. (c) $9 \times 10^{-7}$ par hertz (un par
mégahertz) à \qty{1}{GHz}~; $2 \times 10^5$ par hertz à \qty{500}{THz}. (d) Les modes sont
discrets quand leur espacement dépasse leur largeur~; dans le four, un mode
par \qty{10}{MHz} environ autour de \qty{2.45}{GHz} --- des dizaines à quelques
pour cent.
\end{solution}

\begin{solution}{exo:b2:guided-waves:10}
(a) De $\operatorname{\vect{curl}}\vect E = -\partial_t\vect B$~:
\[
\underline B_x = -\frac{k_g}\omega\underline E , \qquad
\underline B_z = \frac{\iu\pi}{a\omega}E_0\cos\frac{\pi x}a\,\eu^{\iu(\omega t - k_gz)} ,
\]
et $\partial_xB_x + \partial_zB_z = -(k_g\pi/a\omega)E_0\cos(\pi x/a) + (k_g\pi/a\omega)E_0
\cos(\pi x/a) = 0$. (b) $(\operatorname{\vect{curl}}\vect B)_y = \partial_zB_x - \partial_xB_z =
(\iu/\omega)(k_g^2 + \pi^2/a^2)\underline E = (\iu\omega/c^2)\underline E$~: $k_g^2 = \omega^2/c^2 - \pi^2/a^2$.
(c) Grandes faces~: $\vect j_s = \pm(B_z\vect e_x - B_x\vect e_z)/\mu_0$ --- selon $z$
en $\sin(\pi x/a)$, en travers en $\cos(\pi x/a)$ (nul au milieu)~; petites
faces~: selon $y$ seulement. (d) Une fente longitudinale centrale ne coupe que les
courants transversaux, qui y sont nuls~; une fente transversale coupe les
courants longitudinaux et rayonne~: c'est l'antenne à fentes.
\end{solution}

\begin{solution}{exo:b2:guided-waves:11}
(a) $2a\sqrt{n_1^2 - n_2^2}/\lambda = 2 \times 50 \times 0.242/1.5 \approx 16$ par polarisation.
(b) $a < \lambda/2\mathrm{NA} = \qty{3.1}{\micro m}$. (c) Un cœur circulaire d'$\mathrm{NA} =
0.12$ est monomode en dessous de $d = 2.405\lambda/\pi\mathrm{NA} = \qty{10}{\micro m}$~: d'où le
standard à \qty{9}{\micro m}. (d) Un rayon rasant est toujours totalement réfléchi~: le
mode le plus bas existe à toute fréquence, ses queues évanescentes s'étalant
simplement davantage dans la gaine.
\end{solution}

\begin{solution}{exo:b2:guided-waves:12}
(a) Le rayon hors axe passe son temps dans un indice plus faible, où la lumière est
plus rapide~: le trajet plus long est compensé. (b) $n(r)\cos\theta = n_1\cos\theta_0$
avec $n \approx n_1(1 - \Delta r^2/a^2)$ et $\cos\theta \approx 1 - \theta^2/2$~: $(\dd r/\dd z)^2 =
\theta_0^2 - 2\Delta r^2/a^2$, une oscillation harmonique $r = r_0\sin(\sqrt{2\Delta}\,z/a)$.
(c) $2\pi a/\sqrt{2\Delta} = \qty{1.1}{mm}$. (d) $n_1L\Delta^2/2c = \qty{0.25}{ns/km}$ contre
\qty{49}{ns/km}~: deux cents fois mieux.
\end{solution}

\begin{solution}{pb:b2:guided-waves:1}
\textbf{1.} $c/2a = \qty{1.74}{GHz}$~; TE$_{20}$ et TE$_{01}$~: \qty{3.49}{GHz}~: seul
TE$_{10}$ à \qty{2.45}{GHz}.

\textbf{2.} $\sqrt{1 - (1.74/2.45)^2} = 0.70$~: $\lambda_g = \qty{17.4}{cm}$, $v_\varphi = 1.42c$,
$v_g = 0.70c = \qty{2.1e8}{m/s}$.

\textbf{3.} $E_0^2 = 4P/\varepsilon_0c\,ab\sqrt{\cdot} = 3200/6.9 \times 10^{-6}$~: $E_0 = \qty{21}{kV/m}$,
soit un centième du claquage.

\textbf{4.} $j_s \approx E_0/\mu_0c = \qty{57}{A/m}$~; $\delta = \qty{10}{\micro m}$, $R_s = \qty{0.1}{\ohm}$~;
$\tfrac12R_sj_s^2 \times 2a = \qty{28}{W/m}$~: $3.5\%$ par mètre --- l'acier est dissipatif,
mais le guide est court.

\textbf{5.} $f = 1.5 \times 10^8\sqrt{11.1(m^2 + n^2) + 25p^2}$ à $\pm3\%$ de
\qty{2.45}{GHz}~: $(1, 2, 3)$, $(2, 1, 3)$ à \qty{2.51}{GHz}, $(4, 0, 2)$, $(0, 4, 2)$,
$(4, 3, 0)$, $(3, 4, 0)$ à \qty{2.50}{GHz} --- environ six.

\textbf{6.} Demi-longueurs d'onde $a/m$, $b/n$, $d/p$~: de $7$ à \qty{10}{cm} d'écart. Le
plateau tournant entraîne les aliments à travers maxima et minima~; un brasseur (un
ventilateur métallique tournant) rebat les modes.

\textbf{7.} $f/Q = \qty{12}{MHz}$~: les résonances chargées se recouvrent en un
continuum~; le magnétron trouve toujours un mode à alimenter.

\textbf{8.} $W = QP/\omega = 3000 \times 800/1.54 \times 10^{10} = \qty{0.16}{J}$~; $u = \qty{8.7}{J/m^3}$~;
$E_0 \approx \sqrt{4u/\varepsilon_0} = \qty{2}{MV/m}$ --- près du claquage~: des arcs~; la puissance
réfléchie retourne au magnétron, protégé (un peu) par une charge
fictive ou un circulateur --- d'où le «~ne jamais le faire tourner à vide~».

\textbf{9.} $0.2/2.1 \times 10^8 = \qty{1}{ns}$.

\textbf{10.} \qty{1.6}{kW} crête~; entre les salves le champ décroît en
$Q/\omega = \qty{13}{ns}$~: la cavité est vide la plupart du temps.

\textbf{11.} $\ell \approx a/\pi = \qty{0.48}{mm}$~; $\eu^{-1/0.48}$~: \qty{-18}{dB} en amplitude,
\qty{-36}{dB} en puissance.

\textbf{12.} $\lambda/4 = \qty{3.1}{cm}$.

\textbf{13.} \qty{5}{mW/cm^2} $\times$ \qty{100}{cm^2} = \qty{0.5}{W}~: soit $0.06\%$.

\textbf{14.} $f_1 \approx c/2d = \qty{3.75}{GHz} > \qty{2.45}{GHz}$~: sous la coupure~; $\ell =
\qty{1.7}{cm}$, $\eu^{-30/1.7}$~: \qty{-150}{dB}.

\textbf{15.} Les courants de paroi doivent passer sans rupture dans la grille~;
une rupture de la liaison est une fente qui les coupe --- une antenne à fente
qui rayonne vers l'extérieur.

\textbf{16.} $\mathrm{NA} = \sqrt{1.465^2 - 1.460^2} = 0.12$, \qty{7}{\degree}~: un cône
étroit, qui exige un laser et une lentille.

\textbf{17.} $(n_1L/c)(n_1/n_2 - 1) = \qty{17}{ns/km}$~; \qty{170}{ns} sur \qty{10}{km}~:
\qty{6}{Mbit/s}.

\textbf{18.} $\tfrac12(\pi d\,\mathrm{NA}/\lambda)^2 = 75$~; pour \qty{9}{\micro m}~: $2.4$ --- le
régime monomode ($V = \pi d\,\mathrm{NA}/\lambda = 2.2 < 2.4$).

\textbf{19.} $\Delta x_0 = (c/n)\tau = \qty{1}{cm}$~: $t_{\text{ét}} = 10^{-4}/0.2 = \qty{0.5}{ms}$,
\qty{100}{km}~: \qty{10}{Gbit/s} sur cette distance avant compensation.

\textbf{20.} \qty{20}{dB}, un facteur $100$~: \qty{10}{\micro W}~; $6000/80 = 75$
amplificateurs.

\textbf{21.} $f = \qty{1.9e14}{Hz}$, $hf = \qty{0.8}{eV}$~; $8 \times 10^{15}$ photons par
seconde~; $8 \times 10^5$ par bit à l'entrée, $8 \times 10^3$ après \qty{100}{km}.

\textbf{22.} La somme des deux pertes est minimale au voisinage de \qty{1.55}{\micro m}
(\qty{0.2}{dB/km}).

\textbf{23.} $R = ((1.465 - 1)/2.465)^2 = 3.6\%$~: \qty{0.16}{dB} par face, deux fois
à un connecteur, plus l'interférence de l'intervalle~; un gel ou un contact physique
supprime l'air.

\textbf{24.} Dans une courbe douce le rayon rencontre encore la frontière au-delà
de l'angle critique~; dans une courbe serrée l'incidence sur la paroi extérieure
passe en dessous et la lumière fuit.

\textbf{25.} Four~: parois métalliques, réflexion avec $r = -1$~; bande passante fixée
par les modes de la cavité~; pertes dans la peau de l'acier. Fibre~: \omterm{thm:b2:wave-interfaces:snell}{réflexion
totale} à une frontière verre--verre~; bande passante fixée par les \omterm{prop:b2:guided-waves:fibre}{dispersions modale} et
chromatique~; pertes par diffusion et absorption, \qty{0.2}{dB/km}.
\end{solution}
